
\input texsis
\paper

\overfullrule=0pt
\def\frac#1#2{#1\over #2}
%
\def\bmatrix#1{\left[ \matrix{#1} \right]}
\def\be{{\beta}}
\def\toone{{t+1}}

%\def\frac#1#2{#1\over #2}
%\magnification=1200
\overfullrule=0pt

\line{Advanced Macro and Time Series \hfil SMU}
\line{Winter 2019 \hfil Thomas Sargent}

%\centerline{\bf Practice problems}
\centerline{\bf Take Home Final Examination}


\medskip

\noindent{\bf Instructions:}  Please feel free to work  in groups of 3 or 4 students. But please hand in your own work. This take home final is
due on  Thursday Feb 28.

\medskip

\noindent{\bf Multiple roles of $\theta_t$}

\medskip

\noindent This question is based on the quantecon lecture ``Ramsey Plans, Time Inconsistency, Sustainable Plans'' -- the second lecture in the Python
``Dynamic Programming Squared'' section, which  we studied in detail in last Thursday's class. 

\medskip
\noindent The lecture shows how to  represent a Ramsey plan
$ \vec \mu $ recursively with the following system created in the spirit of Chang in the quantecon lecture. 
In particular, in the end we obtained the recursive representation:
$$
\eqalign{
\theta_0 & = \theta_0^* \cr
\mu_t &  = b_0 + b_1 \theta_t \cr
\theta_{t+1} & = d_0 + d_1 \theta_t } \leqno(1) $$

\medskip

\noindent We can think of the  sequence
$ \{\theta_t\}_{t=0}^\infty $ as a sequence of
synthetic ``promised inflation rates'' that are just computational devices for
generating a sequence $ \vec\mu $ of money growth rates that are to
be substituted into the following equation in order to form actual rates of inflation:

$$
\theta_t = {\frac{1}{1+\alpha}} \sum_{j=0}^\infty \left(\frac{\alpha}{1+\alpha}\right)^j \mu_{t+j}. \quad \leqno(2)
$$

\noindent where $\alpha \in (0,1)$.

\medskip
\noindent{\bf a.} Please describe in detail an algorithm that compute the parameters $\theta_0^*, b_0, b_1, d_0, d_1$ that appear in
system (1).  
\medskip

\medskip
\noindent{\bf b.}  Please tell why the inflation rate $ \theta_t $ that appears in system (1) and 
equation (2) can be thought of as  playing three roles simultaneously:

  * In equation (2), $ \theta_t $ is the actual rate of inflation
  between $ t $ and $ t+1 $
  
  * $ \theta_t $ is also the public's
  expected rate of inflation between $ t $ and $ t+1 $

  * $ \theta_t $ is a promised rate of inflation
  chosen by the Ramsey planner at time $ 0 $  


\medskip 
\noindent{\bf c.} Please provide  representation of the Ramsey planner's optimal  $\vec \mu$, i.e., its Ramsey plan.


\medskip
\noindent{\bf d.} Please verify that your representation  of a Ramsey plan guarantees that promised
inflation equals actual inflation. 


\medskip \noindent {\it Hint:} Some  ``Big $K$, little $k$'' logic might help you.


\bye






\bye
